% Fragmento público generado desde propietarios íntegros.
% Residencia: 52.10.1.
% Fuente: ../incoming/radion_monografia_20260827/manuscrito/sections/11_modularidad_orbitas.tex:1-84
\noindent Les caractères modulaires, orbitaux et de chemin conservent des invariants distincts d'une même forme orientée.\par\medskip

\subsubsection{Tore complexe associé à l'ellipse d'action}

\paragraph{Réseau rectangulaire de périodes}

L'ellipse à elle seule n'est pas un tore. Nous déclarons le réseau
\begin{equation}
\Lambda_\eta=a_S\Z+i,b_S\Z
\label{eq:lattice}
\end{equation}
et le quotient \(\C/\Lambda_\eta\). Son module est
\begin{equation}
\tau_\eta=i\frac{b_S}{a_S}=ie^{-\eta}
=i\,0.741048144159375160795483663369300\ldots.
\label{eq:tau-eta}
\end{equation}
L'inversion bidirectionnelle \(\eta\mapsto-\eta\) échange les axes :
\begin{equation}
\tau_{-\eta}=ie^\eta=-\frac1{\tau_\eta}.
\label{eq:modular-S}
\end{equation}
C'est exactement la transformation modulaire \(S:\tau\mapsto-1/\tau\).

\paragraph{Invariance de la fonction modulaire \(j\)}

Par invariance modulaire,
\begin{equation}
\boxed{j(\tau_{-\eta})=j(\tau_\eta).}
\label{eq:j-invariance}
\end{equation}
Dans la normalisation \(j(i)=1728\),
\[
j(\tau_\eta)
=5597.43843932408729830097089254768\ldots.
\]
La fonction moonshine normalisée \(J=j-744\) donne
\[
J(\tau_\eta)
=4853.43843932408729830097089254768\ldots.
\]
Pour le bloc APP basal,
\[
\tau_0=i\frac{\sqrt5}{3},
\qquad
j(\tau_0)=5369.48832024674306021916882626352\ldots.
\]

\begin{proposition}[Ce que conserve la classe modulaire]
L'involution de feuillet inverse \(\eta\), mais \(j\) conserve la classe du
tore. Les foyers et l'étiquetage des axes retiennent l'orientation que \(j\)
oublie.
\end{proposition}

Ce résultat relie de manière exacte la double voie à la modularité une
fois le quotient déclaré. Il n'autorise pas à faire agir directement le
Monstre sur la charge, les chiffres ou les foyers. Le couloir exceptionnel
\[
\mathrm{APP}\to\mathrm{Paley}\to\mathrm{Witt}\to\mathrm{Golay}
\to A_2^{12}\to\mathrm{Leech}\to V^\natural
\to\mathrm{Monster}\to\mathrm{Moonshine}
\]
est une autre projection du même état enrichi. Pour identifier une action
concrète sur le radion, il faut un entrelaceur qui préserve ses
coordonnées pertinentes.

\paragraph{Chaîne d'invariants du caractère de forme}

En réunissant les réalisations déclarées,
\begin{equation}
\boxed{
r_\eta=e^{-\eta}
=\sqrt{\frac{A-C^*}{A+C^*}}
=\frac{b_1}{a_1}
=\frac{b_S}{a_S}
=\frac{\Delta P}{\Delta Q}
=\frac{Z_\eta}{Z_*}
=\Im\tau_\eta.}
\label{eq:shape-master}
\end{equation}
Un seul rapport parcourt géométrie spectrale, action, covariance, impédance et
module. Chaque égalité repose sur une application explicite ; aucune ne dépend d'une
ressemblance décimale.

